% =====================================================================
\section{Numerical experiments}\label{sec:experiments}
% =====================================================================

We illustrate the theory on a scalar double-well benchmark chosen so
that all hypotheses of \cref{ass:standing} hold and the exact
quasi-static branches are known in closed form.

\subsection{Benchmark}\label{sec:benchmark}

Let $\Q=\R$ and
\begin{equation}\label{eq:benchmark}
  W(q)=k\,q^2(1-q)^2,\qquad
  \ell(S)=\ell_\infty\bigl(1-e^{-\lambda S}\bigr),\qquad
  \Psi(\dot q)=\rho\,|\dot q|,
\end{equation}
with $k,\ell_\infty,\lambda,\rho>0$, so that
$E(q,S)=k\,q^2(1-q)^2-q\,\ell(S)$ and
$g(q)=W'(q)=2k\,q(1-q)(1-2q)$. The loading $\ell$ is strictly
increasing and saturating: a \emph{high} signal $S$ drives $q$ toward
the second well $q=1$. It is bounded ($C_\ell=\ell_\infty$) and
Lipschitz ($C_{\ell'}=\ell_\infty\lambda$), so~\ref{A:ell} holds;
$W$ is a coercive quartic, so~\ref{A:W} holds with $p=4$; and
$\Psi=\rho|\cdot|$ satisfies~\ref{A:Psi} and~\ref{A:coerc} with
$c_\Psi=\rho$. The energy $E(\cdot,S)$ has at most two local minima,
so \cref{thm:BV_uniqueness} guarantees a unique BV solution.

The lower branch of $g$ attains its maximum
\begin{equation}\label{eq:gmax}
  g_{\max}=\frac{k}{3\sqrt3}
  \qquad\text{at}\qquad q_-=\frac{3-\sqrt3}{6},
\end{equation}
the fold beyond which the near branch of $g(q)=\tau$ disappears and the
system jumps. On the yield surface $g(q)=\tau$ (with
$\tau=\ell-\sgn(f)\rho$), the admissible root on $[0,q_-]$ is given by
the trigonometric (Cardano) formula
\begin{equation}\label{eq:cubic}
  q=\tfrac12+\tfrac{1}{\sqrt3}\cos\!\Bigl(\tfrac13\arccos\!\bigl(\tfrac{\tau}{g_{\max}}\bigr)-\tfrac{4\pi}{3}\Bigr),
  \qquad |\tau|<g_{\max}.
\end{equation}
We drive the system with the smooth cyclic loading
$S(t)=\tfrac12 S_{\max}\bigl(1-\cos(2\pi t/T)\bigr)$, $t\in[0,T]$, which
increases from $0$ to $S_{\max}$ and back. We fix $k=1$, $\rho=0.03$,
$\lambda=1$, $S_{\max}=5$, $T=1$, and consider two regimes:
a \emph{reversible} one ($\ell_\infty=0.18$, so
$\ell_\infty-\rho=0.15<g_{\max}\approx0.1925$) and a
\emph{catastrophic} one ($\ell_\infty=0.30$, so
$\ell_\infty-\rho=0.27>g_{\max}$).

\subsection{Qualitative behaviour}\label{sec:qualitative}

\Cref{fig:landscape} shows how the loading tilts the double-well
landscape, lowering the barrier toward $q=1$ as $S$ grows.
\Cref{fig:evolution}(a) shows the reversible cycle: the state remains
locked (elastic) until the driving force reaches the yield threshold,
then flows, and on unloading locks again before partially recovering,
leaving a residual $q_{\mathrm{res}}>0$. \Cref{fig:evolution}(b) shows
the resulting hysteresis loops; in the catastrophic regime the near
branch disappears at the fold and the state jumps to the far well
($q\approx1$), the signature of an abrupt transition.

The closed-form predictions~\eqref{eq:cubic} agree with the computed
extrema to four significant digits, as reported in
\cref{tab:validation}.

\begin{figure}[t]
  \centering
  \begin{tikzpicture}
    \begin{axis}[width=0.62\textwidth,height=0.34\textwidth,
      xlabel=$q$,ylabel=$E(q,S)$,
      legend pos=north west,legend cell align=left,
      grid=both,grid style={gray!20},tick label style={font=\small}]
      \addplot[thick,blue] table[x=q,y=E]{figures/fig_landscape_lo.dat};
      \addplot[thick,teal!70!black] table[x=q,y=E]{figures/fig_landscape_mid.dat};
      \addplot[thick,red!70!black] table[x=q,y=E]{figures/fig_landscape_hi.dat};
      \legend{$S=0$,$S=1$,$S=5$}
    \end{axis}
  \end{tikzpicture}
  \caption{Tilting of the double-well energy $E(\cdot,S)$ as the loading
  $S$ increases (reversible regime, $\ell_\infty=0.18$). The right well
  is progressively favoured.}
  \label{fig:landscape}
\end{figure}

\begin{figure}[t]
  \centering
  \begin{subfigure}[t]{0.49\textwidth}
    \centering
    \begin{tikzpicture}
      \begin{axis}[width=\textwidth,height=0.72\textwidth,
        xmin=0,xmax=1,axis y line*=left,xlabel=$t$,ylabel=$q(t)$,
        ymin=-0.01,tick label style={font=\small}]
        \addplot[thick,blue] table[x=t,y=q]{figures/fig_time_smooth.dat};
      \end{axis}
      \begin{axis}[width=\textwidth,height=0.72\textwidth,
        xmin=0,xmax=1,axis y line*=right,axis x line=none,ylabel=$S(t)$,
        tick label style={font=\small},legend pos=north east,
        legend cell align=left]
        \addlegendimage{thick,blue}\addlegendentry{$q(t)$}
        \addplot[thick,red!70!black,dashed] table[x=t,y=S]{figures/fig_time_smooth.dat};
        \addlegendentry{$S(t)$}
      \end{axis}
    \end{tikzpicture}
    \caption{Time evolution (reversible)}
  \end{subfigure}
  \hfill
  \begin{subfigure}[t]{0.49\textwidth}
    \centering
    \begin{tikzpicture}
      \begin{axis}[width=\textwidth,height=0.72\textwidth,
        xlabel=$\ell(S)$,ylabel=$q$,legend pos=north west,
        legend cell align=left,grid=both,grid style={gray!20},
        tick label style={font=\small}]
        \addplot[thick,blue] table[x=ell,y=q]{figures/fig_loop_smooth.dat};
        \addplot[thick,red!70!black] table[x=ell,y=q]{figures/fig_loop_cat.dat};
        \legend{reversible,catastrophic}
      \end{axis}
    \end{tikzpicture}
    \caption{Hysteresis loops}
  \end{subfigure}
  \caption{(a) Elastic lock, dissipative flow and partial recovery over
  a loading cycle. (b) Hysteresis loops in the $(\ell,q)$ plane: the
  reversible loop stays on the near branch, while the catastrophic loop
  jumps at the fold $g_{\max}$ to the far well.}
  \label{fig:evolution}
\end{figure}

\begin{table}[t]
  \centering
  \caption{Analytical~\eqref{eq:cubic} vs.\ computed extrema
  (reversible regime).}
  \label{tab:validation}
  \begin{tabular}{lccc}
    \toprule
    & analytical & numerical & rel.\ error \\
    \midrule
    $q_{\max}$ (peak damage)   & $0.10540$ & $0.10535$ & $5\cdot10^{-4}$ \\
    $q_{\mathrm{res}}$ (residual scar) & $0.01578$ & $0.01573$ & $3\cdot10^{-3}$ \\
    \bottomrule
  \end{tabular}
\end{table}

\subsection{Convergence and consistency}\label{sec:conv_exp}

We verify the consistency theory of \cref{sec:consistency} on the
reversible cycle. Two error measures are computed against a fine
reference ($N=2\cdot10^5$): the $L^\infty$ error of the state
$\|q^h-q_{\mathrm{ref}}\|_\infty$, and the global energy-balance defect
$\sum_k e_k=(E_0+\mathcal W_h)-(E_N+\Diss_h)\ge0$, the accumulated gap
in the discrete energy inequality~\eqref{eq:disc_energy}.

Two facts emerge, both consistent with the theory. First, the state
error is at the level of machine precision ($\sim10^{-16}$): during
smooth flow the return map~\eqref{eq:scalar_return_map} solves the
algebraic yield condition $g(q_{k+1})=\ell(S_{k+1})-\sgn(f)\rho$
\emph{exactly} at each node, independently of $h$, so the incremental
scheme reproduces the exact quasi-static branch at the grid points.
The entire discretisation error is thus carried by the energy balance,
not by the state. Second, the energy defect decays as $O(h)$
(\cref{fig:convergence}), exactly the first-order global energy
consistency of \cref{thm:global_consistency}: halving $h$ halves the
defect, and the data track the reference slope-one line.

\begin{figure}[t]
  \centering
  \begin{tikzpicture}
    \begin{loglogaxis}[width=0.6\textwidth,height=0.4\textwidth,
      xlabel=$h$,ylabel=$\sum_k e_k$ (energy defect),
      legend pos=north west,legend cell align=left,
      grid=both,grid style={gray!20},tick label style={font=\small}]
      \addplot[only marks,mark=*,blue] table[x=h,y=energy_defect]{figures/fig_convergence.dat};
      \addplot[dashed,black,domain=1.5e-4:2e-2]{0.0748*x};
      \legend{energy defect,slope $1$}
    \end{loglogaxis}
  \end{tikzpicture}
  \caption{Global energy-balance defect versus step size $h$
  (log--log). The first-order decay confirms
  \cref{thm:global_consistency}. The state error, by contrast, is at
  machine precision because the return map is exact on smooth branches.}
  \label{fig:convergence}
\end{figure}
